\section{From the Mountain to a Universal Memorial}

\subsection{Hermits, an oratory, and a rule}

Latin Christian hermits gathered near the spring on Mount Carmel around the turn of the thirteenth century. Their own first surviving normative witness is the \textit{formula vitae} given at their request by Albert, Latin Patriarch of Jerusalem, between 1206 and 1214. It addresses ``B. and the other hermits'' and says nothing of a founding apparition to Elijah or Simon Stock. Its first positive program is allegiance to Jesus Christ.

The hermits built an oratory among their cells, as the Rule required, and dedicated it to the Blessed Virgin. In the patronage language of the period they became the Brothers of the Blessed Virgin Mary of Mount Carmel. Mary was the Lady of their place and the patroness under whose protection they lived; this Marian identity precedes the surviving Scapular-vision narrative.

Honorius III's \textit{Ut vivendi normam} of 30 January 1226 commanded observance of Albert's pre-Lateran form of life. Gregory IX confirmed it in 1229. After migration toward Europe and a changed mendicant setting, Innocent IV's \textit{Quae honorem Conditoris} of 1 October 1247 approved adaptations, including provisions suited to common refectory and urban life. The Order did not abandon contemplation when it became mendicant; it learned to hold cell, choir, common life, preaching, and service in one vocation.

\subsection{A historical spine}

\begin{center}
\footnotesize
\begin{longtable}{@{}p{0.16\linewidth}p{0.34\linewidth}p{0.43\linewidth}@{}}
\toprule
\textbf{Date} & \textbf{Development} & \textbf{What it establishes---and does not}\\
\midrule
\endhead
ca.~1206--1214 & Albert gives the hermits a \textit{formula vitae}. & Establishes Christ-centered Carmelite life near the spring; not an Elijah-foundation charter or Scapular account.\\
30 Jan.~1226 & Honorius III, \textit{Ut vivendi normam}. & Papal confirmation of the received form of life.\\
1 Oct.~1247 & Innocent IV, \textit{Quae honorem Conditoris}. & Definitive adapted Rule for changed conditions in the West.\\
1274--1298 & Lyons II threatens newer orders; Carmelite continuance is preserved and later secured. & Helps explain the Order's strong memory of Marian patronage and deliverance; does not by itself prove a 1251 apparition.\\
late 14th c. & A written Simon-vision tradition becomes visible; the Carmelite Marian commemoration develops. & Documents reception significantly later than the alleged event.\\
15th--16th c. & Scapular confraternities and the vision narrative expand; the \textit{Flos Carmeli} is received in Carmelite worship. & Shows devotional growth, not apostolic or thirteenth-century textual continuity for every promise.\\
24 Sept.~1726 & Benedict XIII extends the feast to the Latin Church. & Establishes universal liturgical reception of the title, not a judgment authenticating each origin narrative.\\
11 Feb.~1950 & Pius XII, \textit{Neminem profecto latet}. & Commends the Scapular as a sign of Marian dedication and Christian life.\\
5 Jan.~1996 & CDWDS approves the Rite of Blessing and Enrollment. & Establishes the Church's present ritual object: habit-sign, baptismal remembrance, affiliation, and commitment.\\
25 Mar.~2001 & John Paul II addresses both Carmelite branches. & Receives the 1251 account as venerable Order tradition while grounding the Scapular in conformity to Christ, prayer, sacraments, and mercy.\\
\bottomrule
\end{longtable}
\end{center}

\subsection{Migration, reform, and the many rooms of Carmel}

Political and military instability in the Holy Land pressed the hermits westward during the thirteenth century. European conditions required clarification of identity: neither monks of an imagined uninterrupted Elijah succession nor a new association without roots, they were a papally confirmed mendicant order bearing the place-name and patronage of Carmel.

Across later centuries the family diversified. The ancient observance, enclosed nuns, Third Order and confraternities, Teresa of Jesus's sixteenth-century reform, John of the Cross's collaboration, Discalced friars and nuns, secular branches, affiliated institutes, and other communities share the Rule and Carmelite inheritance in distinct juridic forms. A wearer of the small Scapular is not thereby a friar, nun, or secular Carmelite. Affiliation admits degrees; vocation and promises require their proper discernment and law.

The saints keep the devotion from shrinking into one object. Teresa describes prayer as friendship with the One who loves us. John of the Cross makes the mountain a demanding image of purification and union with God. Thérèse of Lisieux finds a ``little way'' of trusting love. Elizabeth of the Trinity dwells in baptismal Trinitarian life. Titus Brandsma joins contemplative truth to resistance against dehumanizing ideology. Teresa Benedicta of the Cross pursues truth through philosophy, Carmel, and martyrdom. Their differences show the charism's range; none licenses bypassing the Church, neighbor, reason, or the Cross.

\subsection{How 16 July became the feast}

The Order's commemoration of Mary developed in the later fourteenth century and came to be associated with its survival and patronage. Carmelite historical accounts connect the choice of 16 July with a desired commemoration near 17 July, the date associated with deliverance at Lyons but already occupied by Saint Alexius. As the Simon Stock story spread, 16 July also became the traditional date of the Scapular gift.

By the late fifteenth century the Carmelite family kept the day as its chief Marian celebration. The feast spread by territorial grants in the seventeenth century and was extended to the Latin Church by Benedict XIII on 24 September 1726. The postconciliar General Roman Calendar retains 16 July as an optional memorial; Carmelite proper calendars celebrate their patroness with greater solemnity. These layers explain why ``feast,'' ``memorial,'' and ``solemnity'' may all appear in reliable sources: they refer to different calendars, places, or periods.

\subsection{The history as spiritual correction}

The devotion does not become less Catholic when its layers are named. The earliest secure core is stronger than a slogan: a community living under a Christ-centered Rule, gathered at the Eucharist, pondering Scripture, and entrusted to Mary. The feast, Scapular, confraternities, poems, papal commendations, and popular customs then form a genuine history of reception. Historical honesty frees the petitioner to love what the Church actually blesses without making faith depend on an unproved date or sentence.
